% PAGES: 134-135
% Continues Chapter 4, section 4.1 "Definitions and properties" after the
% c1 agent's assignment (pages 127-133). Page 134 (folio 118) opens with a
% fresh "Example:" (the Jordan block matrix, footnoted "A matrix of the
% form (4.9) is called a Jordan block."). This is a new, self-contained
% example -- NOT a continuation of the "Examples:" (plural) block c1 was
% transcribing (matrices A and B revisited via characteristic polynomials),
% which c1's own sec04_c1_2.tex notes closes within c1's page range with a
% manual "\hfill$\square$\par\medskip" after the matrix-B case analysis.
% No environment is left open across the c1/c2 boundary; nothing needed to
% be closed here.
%
% The source prints this example as two separate italic run-in labels,
% "Example:" (statement -- no closing mark) immediately followed by
% "Answer:" (solution -- closes with a small square). Since \examplex
% (preamble.tex) always appends a closing square at \end{examplex}, and no
% square appears after the "Example:" statement on the page, the
% "Example:" label is written out manually here (the same convention c1
% used for the plural "Examples:" label in sec04_c1_2.tex) instead of
% \begin{examplex}. The \answerx environment (preamble.tex) is used for
% "Answer:"; \answerx does NOT auto-append a closing square, so one is
% added manually just before \end{answerx}, matching the square printed on
% the page. This is unrelated to the \begin{proof} auto-qed rule, which
% applies only to the amsthm "proof" environment.

\par\medskip\noindent\textit{Example:}\ The $n \times n$ matrix\footnote{A matrix of
the form (4.9) is called a Jordan block.}
\begin{equation}
A \;=\; \begin{pmatrix}
d & 1 & 0 & \dots & 0\\
0 & d & 1 & \ddots & \vdots\\
\vdots & \ddots & \ddots & \ddots & 0\\
\vdots & \ddots & \ddots & \ddots & 1\\
0 & \dots & \dots & 0 & d
\end{pmatrix}
\end{equation}
has eigenvalue $d$ of multiplicity $n$, and $\begin{pmatrix}1\\0\\\vdots\\0\end{pmatrix}$
is the only eigenvector corresponding to $d$.

\begin{answerx}
From (4.8) and (10.5), it follows that the characteristic polynomial of the upper
triangular matrix $A$ is
\[
P_A(t) \;=\; \det(tI-A) \;=\; (t-d)^n.
\]
Thus, the polynomial $P_A(t)$ has root $d$ with multiplicity $n$, and we conclude that
the matrix $A$ has eigenvalue $d$ with multiplicity $n$.

Let $v = (v_i)_{i=1:n}$ be an eigenvector of the matrix $A$ corresponding to the
eigenvalue $d$. Then, $Av = dv$, i.e.,
\[
\begin{cases}
dv_i + v_{i+1} \;=\; dv_i, & \forall\, i = 1:(n-1);\\
v_n \;=\; 0
\end{cases}
\]
\[
\Longleftrightarrow\quad
\begin{cases}
v_{i+1} \;=\; 0, & \forall\, i = 1:(n-1);\\
v_n \;=\; 0
\end{cases}
\]
Thus, $v = \begin{pmatrix}v_1\\0\\\vdots\\0\end{pmatrix} =
v_1\begin{pmatrix}1\\0\\\vdots\\0\end{pmatrix}$, and we conclude that
$\begin{pmatrix}1\\0\\\vdots\\0\end{pmatrix}$ is the only eigenvector corresponding to
the eigenvalue $d$ of the matrix $A$. \hfill$\square$
\end{answerx}

The number of eigenvectors of a matrix is, by definition, equal to the number of
linearly independent eigenvectors of the matrix. It is important to note that
eigenvectors corresponding to different eigenvalues are linearly independent; see
Lemma 4.2 below:

\begin{lemma}
The eigenvectors corresponding to different eigenvalues of a matrix are linearly
independent.
\end{lemma}

\begin{proof}
We give a proof by contradiction. Let $A$ be a square matrix of size $n$, and assume
that there exist eigenvectors of $A$, each one corresponding to a different eigenvalue
of $A$, which are not linearly independent. Since the matrix $A$ has a finite number
(not larger than $n$) of different eigenvalues, there must exist a \emph{minimal}
number of linearly dependent eigenvectors of $A$ with each vector corresponding to a
different eigenvalue. Denote this minimal number by $p$.

Let $v_1, v_2, \dots, v_p$ be $p$ linearly dependent eigenvectors of $A$ corresponding
to $p$ distinct eigenvalues of $A$ denoted be $\lambda_1, \lambda_2, \dots, \lambda_p$,
i.e., such that
\begin{equation}
c_1 v_1 \;+\; c_2 v_2 \;+\; \dots \;+\; c_p v_p \;=\; 0,
\end{equation}
with $c_1, c_2, \dots, c_p$ constants such that $c_i \neq 0$, for all $i = 1:p$. By
multiplying (4.10) by the matrix $A$ to the left, we obtain that
\begin{equation}
c_1 A v_1 \;+\; c_2 A v_2 \;+\; \dots \;+\; c_p A v_p \;=\; 0.
\end{equation}
Note that
\begin{equation}
A v_i \;=\; \lambda_i v_i, \quad \forall\, i = 1:p,
\end{equation}
since $v_i$ is an eigenvector corresponding to the eigenvalue $\lambda_i$ of $A$.

From (4.11) and (4.12), it follows that
\begin{equation}
c_1 \lambda_1 v_1 \;+\; c_2 \lambda_2 v_2 \;+\; \dots \;+\; c_p \lambda_p v_p \;=\; 0.
\end{equation}

By multiplying (4.10) by $\lambda_1$, we find that
\begin{equation}
c_1 \lambda_1 v_1 \;+\; c_2 \lambda_1 v_2 \;+\; \dots \;+\; c_p \lambda_1 v_p \;=\; 0,
\end{equation}
and, by subtracting (4.14) from (4.13), we obtain that
\begin{equation}
c_2(\lambda_2-\lambda_1) v_2 \;+\; \dots \;+\; c_p(\lambda_p-\lambda_1) v_p \;=\; 0.
\end{equation}

The coefficients of the vectors $v_2, \dots, v_p$ from (4.15) are nonzero, since
$c_i \neq 0$, for all $i = 2:p$, and $\lambda_i \neq \lambda_1$, for all $i = 2:p$,
since the eigenvalues $\lambda_1, \lambda_2, \dots, \lambda_p$ are distinct.

Thus, we found $p-1$ eigenvectors of the matrix $A$, i.e., $v_2, \dots, v_p$, each one
corresponding to a different eigenvalue of $A$, which are linearly dependent. This
contradicts the fact that $p$ denotes the \emph{minimal number} of linearly dependent
eigenvectors of $A$.

We conclude that eigenvectors corresponding to different eigenvalues of a matrix must
be linearly independent.
\end{proof}

\par\medskip\noindent\textit{Summary: General properties of eigenvalues and
eigenvectors}\par\medskip

\begin{itemize}
\item A square matrix of size $n$ has exactly $n$ eigenvalues (possibly complex
numbers), if the eigenvalues are counted with their multiplicities.

\item A square matrix of size $n$ has at most $n$ linearly independent eigenvectors,
but may have less than $n$ linearly independent eigenvectors.

\item There is at least one eigenvector and there are at most $m_\lambda$ linearly
independent eigenvectors (and possibly fewer) corresponding to an eigenvalue $\lambda$
with multiplicity $m_\lambda$.
\end{itemize}

Note that, while the eigenvalues of a matrix with real entries may be complex numbers,
all the eigenvalues of a symmetric matrix with real entries are real numbers.
Moreover, symmetric matrices have a full set of eigenvectors; see section 5.1 for
details.
